% =====================================================================================
%  RAPPORT TECHNIQUE — AI Project Reviewer
%  Squelette : chaque section indique QUI la rédige (M1..M7, voir GUIDE_PROJET.md) et CE QU'ELLE
%  DOIT CONTENIR. Rédigez avec vos propres mots : le jury vérifie que vous maîtrisez vos choix.
%  Compiler : pdflatex rapport.tex (deux fois)
% =====================================================================================
\documentclass[11pt,a4paper]{report}
\usepackage[utf8]{inputenc}
\usepackage[T1]{fontenc}
\IfFileExists{lmodern.sty}{\usepackage{lmodern}}{}
\IfFileExists{french.ldf}{\usepackage[french]{babel}}{}
\usepackage[margin=2.3cm]{geometry}
\usepackage{booktabs,tabularx,longtable,xcolor,hyperref,listings}
\lstset{basicstyle=\ttfamily\small,breaklines=true,frame=single,language=Java}
\newcommand{\todo}[1]{\par\noindent\colorbox{yellow!30}{\parbox{\dimexpr\linewidth-2\fboxsep}{\textbf{À rédiger :} #1}}\par}

\title{AI Project Reviewer\\\large Plateforme d'évaluation automatique de projets assistée par IA}
\author{Groupe de 7 : M1 \and M2 \and M3 \and M4 \and M5 \and M6 \and M7}
\date{\today}

\begin{document}
\maketitle
\tableofcontents

\chapter{Problématique}                                  % M1
\todo{Pourquoi évaluer automatiquement ? Pourquoi un LLM seul ne suffit pas (non déterministe, contexte limité, manipulable) ? Les 8 problèmes d'architecture listés au §1 du sujet.}

\chapter{Exigences fonctionnelles}                       % M1
\todo{Les 10 fonctionnalités du §2 du sujet, sous forme de tableau « exigence / classe(s) qui la réalisent ».}

\chapter{Architecture générale}                          % M1
\todo{Couches (ui, application, analysis, llm, project, security, report, persistence, config). Règle de dépendance : l'IHM ne connaît que ReviewerFacade ; seul AppFactory connaît les classes concrètes.}

\section{Diagramme des composants principaux}            % M1
\todo{Diagramme UML de composants ou de packages (PlantUML/draw.io exporté en PDF, puis \textbackslash includegraphics).}

\section{Parcours d'une analyse}                          % M5
\todo{Diagramme de séquence : clic « Lancer » → Facade → EvaluationEngine → Analyzer → LlmProvider (décorateurs) → parser → agrégation → Observer → rapport.}

\chapter{Choix technologiques}                           % M1
\todo{Java 17 sans framework (architecture visible), Swing (inclus au JDK, aucune dépendance), Jackson (JSON), JUnit 5, API HTTP « compatible OpenAI » (couvre LM Studio, Ollama, Mistral, DeepSeek), Docker, pdflatex. Pour chaque choix : alternative écartée et pourquoi.}

\chapter{Répartition des tâches}                         % M1
\todo{Tableau membre / responsabilité / packages / tests / sections du rapport.}

\chapter{Design Patterns}                                % chacun rédige SES patterns
\todo{Pour CHAQUE pattern : problème concret rencontré, solution, classes concernées (extrait de code court), bénéfice mesurable (ex. « ajouter un fournisseur = 0 ligne modifiée »), limite. Ajouter une section « patterns volontairement non utilisés ».}
\section{Strategy}          % M5 (Analyzer) et M2 (FileSelectionStrategy)
\section{Adapter}           % M3
\section{Decorator}         % M3
\section{Factory}           % M3 (LlmProviderFactory), M2 (ProjectLoaderFactory), M5 (AnalyzerRegistry)
\section{Template Method}   % M5
\section{Builder}           % M4 (PromptBuilder), M7 (LatexDocumentBuilder)
\section{Composite}         % M2
\section{Observer}          % M5 / M7
\section{Facade}            % M1
\section{Patterns non retenus} % M1

\chapter{Fonctionnement du moteur d'analyse}              % M5
\todo{Registre d'analyseurs, reprise partielle (un critère en échec n'arrête pas les autres), calcul de la note globale, analyses déterministes vs LLM, faits mesurés injectés dans les prompts.}

\chapter{Intégration des LLM}                             % M3
\todo{Interface LlmProvider, adaptateur, empilement des décorateurs, choix du modèle local, justification.}

\section{Gestion du contexte}                             % M4
\todo{Sélection, découpage (Chunker), résumé structurel (StructureSummarizer), agrégation, limite du nombre d'appels, cache.}

\chapter{Stratégie de construction des prompts}           % M4
\todo{Rôle, critère, faits, données encadrées par nonce, format JSON imposé, rappel final, requête de correction.}

\chapter{Gestion des erreurs}                             % M3 et M4
\todo{Tableau : panne (timeout, HTTP, serveur arrêté, vide, malformé, hors schéma, incohérent) → détection → réaction (retry, repli, correction, statut PARTIAL/FAILED) → test qui le prouve.}

\chapter{Génération du rapport}                           % M7
\todo{Séparation résultat / modèle Java / rendu LaTeX ; échappement (et le bug des crochets trouvé en relisant le PDF) ; structure indépendante du LLM ; compilation PDF sécurisée.}

\chapter{Stratégie de test}                               % M1 (coordination), chacun ses tests
\todo{Pyramide de tests, doublures (ScriptedLlmProvider, faux serveur HTTP, faux ProcessRunner), tests d'intégration sans IHM, résultats (nombre de tests).}

\chapter{Architecture de sécurité}                        % M6
\todo{Modèle de menaces : tableau menace / impact / parade / limite (Zip Slip, zip bomb, liens symboliques, exécution de code, injection de prompt, injection LaTeX, fuite de secrets, épuisement de ressources).}
\section{Isolation Docker}                                % M6
\section{Utilisation de la machine virtuelle}             % M6
\section{Injection de prompt}                             % M4

\chapter{Limites de la solution}                          % tous
\chapter{Améliorations possibles}                         % tous
\chapter{Bilan du projet}                                 % tous (bilan individuel court de chaque membre)

\appendix
\chapter{Réponses aux questions du sujet}                 % M1 relit, chacun rédige la sienne
\begin{enumerate}
  \item Quelles sont les principales responsabilités du système ?  % M1
  \item Comment les avez-vous séparées ?                           % M1
  \item Quels Design Patterns et pourquoi ?                        % M5
  \item Comment remplacer le LLM actuel ?                          % M3
  \item Comment ajouter un critère d'évaluation ?                  % M5
  \item Comment tester sans appeler réellement un LLM ?            % M4
  \item Comment gérer une réponse invalide du modèle ?             % M4
  \item Comment empêcher un projet analysé d'endommager la machine ? % M6
  \item Quels sont les principaux points de couplage ?             % M1
\end{enumerate}
\end{document}
